\documentclass{article}
\usepackage{amsmath,amssymb,amsthm}
\usepackage{algorithm,algorithmic}
\usepackage{hyperref}
\usepackage{enumitem}
\newtheorem{theorem}{Theorem}
\newtheorem{lemma}{Lemma}
\newtheorem{assumption}{Assumption}
\newtheorem{corollary}{Corollary}
\newcommand{\prox}{\operatorname{prox}}
\newcommand{\R}{\mathbb{R}}

\begin{document}

\section*{Proximal Gradient Method (PGM)}

\section*{Mathematical Formulation}
Let $f:\R^{d}\to\R$ be convex and $L$‑smooth (i.e. $\nabla f$ is $L$‑Lipschitz) and let $g:\R^{d}\to\R\cup\{+\infty\}$ be proper, closed and convex (possibly non‑smooth).  
Define the composite objective
\[
F(x)\;:=\;f(x)+g(x),\qquad x\in\R^{d}.
\]

Given a stepsize $\eta>0$, the \emph{proximal gradient update} reads
\begin{equation}
\boxed{\;
x_{k+1}\;=\;\prox_{\eta g}\!\bigl(x_{k}-\eta\,\nabla f(x_{k})\bigr),\qquad k=0,1,2,\dots
\;}
\tag{1}
\end{equation}
where the proximal operator of $g$ is
\[
\prox_{\eta g}(v)\;:=\;\arg\min_{u\in\R^{d}}\Bigl\{ g(u)+\frac{1}{2\eta}\|u-v\|^{2}\Bigr\}.
\]

\section*{Pseudocode}
\begin{algorithm}[H]
\caption{Proximal Gradient Method}
\begin{algorithmic}[1]
\STATE \textbf{Input:} stepsize $\eta>0$, initial point $x_{0}\in\R^{d}$, max iterations $K$.
\FOR{$k=0$ \TO $K-1$}
    \STATE Compute gradient $g_{k}:=\nabla f(x_{k})$.
    \STATE Take a proximal step 
    \[
    x_{k+1}\gets \prox_{\eta g}\bigl(x_{k}-\eta g_{k}\bigr).
    \]
\ENDFOR
\STATE \textbf{Output:} $x_{K}$ (or the best iterate in terms of $F$).
\end{algorithmic}
\end{algorithm}

\section*{Dry Run Example}
Consider the scalar problem
\[
f(w)=(w-3)^{2},\qquad g\equiv0,
\]
so $F(w)=f(w)$.  The proximal operator of $g$ is the identity: $\prox_{\eta g}(v)=v$.

\begin{center}
\begin{tabular}{c|c|c|c}
iteration $k$ & $w_{k}$ & $\nabla f(w_{k})$ & $w_{k+1}=w_{k}-\eta\nabla f(w_{k})$\\ \hline
0 & $0$ & $2(0-3)=-6$ & $0-0.1(-6)=0.6$\\
1 & $0.6$ & $2(0.6-3)=-4.8$ & $0.6-0.1(-4.8)=1.08$\\
2 & $1.08$ & $2(1.08-3)=-3.84$ & $1.08-0.1(-3.84)=1.464$\\
3 & $1.464$ & $2(1.464-3)=-3.072$ & $1.464-0.1(-3.072)=1.7712$
\end{tabular}
\end{center}
The corresponding objective values are $F(w_{k})=(w_{k}-3)^{2}$:
\[
F(w_{0})=9,\;F(w_{1})=5.76,\;F(w_{2})=3.6864,\;F(w_{3})=2.3620.
\]
The values decrease, illustrating the descent property of (1).

\newpage

\section*{Assumptions}
\begin{assumption}[Convexity and Smoothness of $f$]\label{ass:convex-smooth}
\begin{enumerate}[label=(\alph*)]
\item $f$ is convex on $\R^{d}$.
\item $\nabla f$ is $L$‑Lipschitz, i.e.
\[
\|\nabla f(x)-\nabla f(y)\|\le L\|x-y\|,\qquad\forall x,y\in\R^{d}.
\]
\end{enumerate}
\end{assumption}
\begin{assumption}[Convexity of $g$]\label{ass:g}
$g$ is proper, closed and convex (possibly non‑smooth).
\end{assumption}
\begin{assumption}[Stepsize]\label{ass:stepsize}
The stepsize satisfies $0<\eta\le\frac{1}{L}$.
\end{assumption}

\section*{Main Convergence Theorem}
\begin{theorem}\label{thm:rate}
Let $\{x_{k}\}$ be generated by \eqref{eq:1} under Assumptions~\ref{ass:convex-smooth}–\ref{ass:stepsize}.  
For any optimal solution $x^{\star}\in\arg\min_{x}F(x)$ we have for all $K\ge1$
\[
F(x_{K})-F(x^{\star})\;\le\;\frac{\|x_{0}-x^{\star}\|^{2}}{2\eta K}.
\tag{2}
\]
Consequently $F(x_{K})\to F(x^{\star})$ at the rate $O(1/K)$.
\end{theorem}

\section*{Theoretical Properties}
\begin{itemize}
\item \textbf{Stability.}  The non‑expansiveness of the proximal operator (Lemma~\ref{lem:nonexp}) together with the $L$‑smoothness of $f$ guarantees that the iterates remain bounded whenever the level set $\{x:F(x)\le F(x_{0})\}$ is bounded.
\item \textbf{Hyperparameter Sensitivity.}  The condition $\eta\le 1/L$ is both sufficient and (up to a constant) necessary for the $O(1/K)$ guarantee.  Larger stepsizes can lead to the extra term $(L-\frac{1}{\eta})\|x_{k+1}-x_{k}\|^{2}$ becoming positive and destroying monotonic descent.
\item \textbf{Comparison to Baselines.}  When $g\equiv0$, (1) reduces to the classical gradient descent method and Theorem~\ref{thm:rate} recovers the standard $O(1/K)$ rate for smooth convex minimization.  When $f\equiv0$, (1) becomes the proximal point algorithm, which also enjoys $O(1/K)$ convergence under the same stepsize rule.
\end{itemize}

\section*{Discussion}
The proof hinges on two elementary facts:
\begin{enumerate}
\item The \emph{descent lemma} for $L$‑smooth $f$,
\item The \emph{firm non‑expansiveness} of $\prox_{\eta g}$.
\end{enumerate}
Both are invoked in a telescoping argument that eliminates the squared‑norm terms and yields the clean $1/K$ bound.  The result extends verbatim to the stochastic setting provided unbiased gradient estimates with bounded variance are used; the deterministic proof presented below remains the core of any such extension.

\newpage

\section*{Preliminaries (Definitions and Lemmas)}
\begin{definition}[Proximal Operator]\label{def:prox}
For $\eta>0$ and convex $g$, 
\[
\prox_{\eta g}(v)=\arg\min_{u}\Bigl\{g(u)+\frac{1}{2\eta}\|u-v\|^{2}\Bigr\}.
\]
\end{definition}
\begin{lemma}[Firm Non‑Expansiveness]\label{lem:nonexp}
For any $u,v\in\R^{d}$,
\[
\bigl\|\prox_{\eta g}(u)-\prox_{\eta g}(v)\bigr\|^{2}
\le \langle u-v,\;\prox_{\eta g}(u)-\prox_{\eta g}(v)\rangle .
\tag{3}
\]
\end{lemma}
\begin{proof}
Let $p=\prox_{\eta g}(u)$ and $q=\prox_{\eta g}(v)$.  
Optimality conditions give
\[
\frac{1}{\eta}(u-p)\in\partial g(p),\qquad
\frac{1}{\eta}(v-q)\in\partial g(q).
\]
By monotonicity of $\partial g$,
\[
\Big\langle\frac{1}{\eta}(u-p)-\frac{1}{\eta}(v-q),\,p-q\Big\rangle\ge0,
\]
which after rearranging yields \eqref{eq:3}.  ∎
\end{proof}
\begin{lemma}[Descent Lemma for $L$‑smooth $f$]\label{lem:descent}
For any $x,y\in\R^{d}$,
\[
f(y)\le f(x)+\langle\nabla f(x),y-x\rangle+\frac{L}{2}\|y-x\|^{2}.
\tag{4}
\]
\end{lemma}
\begin{proof}
Standard consequence of Lipschitz continuity of $\nabla f$.  ∎
\end{proof}

\section*{Main Proof (Step‑by‑Step Derivation)}
Let $x^{\star}\in\arg\min F$ be arbitrary but fixed.  Define the update \eqref{eq:1}
\[
x_{k+1}= \prox_{\eta g}\bigl(x_{k}-\eta\nabla f(x_{k})\bigr).
\]
Denote the intermediate point
\[
y_{k}=x_{k}-\eta\nabla f(x_{k}).
\tag{5}
\]

\noindent\textbf{[Step 1]} Apply Lemma~\ref{lem:nonexp} with $u=y_{k}$ and $v=x^{\star}$:
\[
\|x_{k+1}-x^{\star}\|^{2}
\le \langle y_{k}-x^{\star},\,x_{k+1}-x^{\star}\rangle .
\tag{6}
\]

\noindent\textbf{[Step 2]} Expand $y_{k}$ using \eqref{eq:5}:
\[
\langle y_{k}-x^{\star},\,x_{k+1}-x^{\star}\rangle
=\langle x_{k}-x^{\star},\,x_{k+1}-x^{\star}\rangle
-\eta\langle\nabla f(x_{k}),\,x_{k+1}-x^{\star}\rangle .
\tag{7}
\]

\noindent\textbf{[Step 3]} Rewrite the first inner product via the identity
\[
2\langle a,b\rangle=\|a\|^{2}+\|b\|^{2}-\|a-b\|^{2}.
\]
With $a=x_{k}-x^{\star}$ and $b=x_{k+1}-x^{\star}$ we obtain
\[
2\langle x_{k}-x^{\star},\,x_{k+1}-x^{\star}\rangle
= \|x_{k}-x^{\star}\|^{2}+\|x_{k+1}-x^{\star}\|^{2}
-\|x_{k+1}-x_{k}\|^{2}.
\tag{8}
\]

\noindent\textbf{[Step 4]} Substitute \eqref{eq:8} into \eqref{eq:7} and then into \eqref{eq:6}:
\[
\|x_{k+1}-x^{\star}\|^{2}
\le\frac12\Bigl(\|x_{k}-x^{\star}\|^{2}+\|x_{k+1}-x^{\star}\|^{2}
-\|x_{k+1}-x_{k}\|^{2}\Bigr)
-\eta\langle\nabla f(x_{k}),\,x_{k+1}-x^{\star}\rangle .
\tag{9}
\]

\noindent\textbf{[Step 5]} Cancel the term $\frac12\|x_{k+1}-x^{\star}\|^{2}$ from both sides of \eqref{eq:9} and multiply by $2$:
\[
\|x_{k+1}-x^{\star}\|^{2}
\le \|x_{k}-x^{\star}\|^{2}
-\|x_{k+1}-x_{k}\|^{2}
-2\eta\langle\nabla f(x_{k}),\,x_{k+1}-x^{\star}\rangle .
\tag{10}
\]

\noindent\textbf{[Step 6]} Add and subtract $f(x_{k})$ and $f(x^{\star})$ using Lemma~\ref{lem:descent} with $x=x_{k}$ and $y=x_{k+1}$:
\[
f(x_{k+1})\le f(x_{k})+\langle\nabla f(x_{k}),\,x_{k+1}-x_{k}\rangle
+\frac{L}{2}\|x_{k+1}-x_{k}\|^{2}.
\tag{11}
\]

\noindent\textbf{[Step 7]} Rearrange \eqref{eq:11} to isolate the inner product:
\[
\langle\nabla f(x_{k}),\,x_{k+1}-x_{k}\rangle
\ge f(x_{k+1})-f(x_{k})-\frac{L}{2}\|x_{k+1}-x_{k}\|^{2}.
\tag{12}
\]

\noindent\textbf{[Step 8]} Replace $\langle\nabla f(x_{k}),\,x_{k+1}-x^{\star}\rangle$ in \eqref{eq:10} by
\[
\langle\nabla f(x_{k}),\,x_{k+1}-x^{\star}\rangle
= \langle\nabla f(x_{k}),\,x_{k+1}-x_{k}\rangle
+ \langle\nabla f(x_{k}),\,x_{k}-x^{\star}\rangle .
\tag{13}
\]

\noindent\textbf{[Step 9]} Insert \eqref{eq:12} into \eqref{eq:13} and then the result into \eqref{eq:10}:
\[
\begin{aligned}
\|x_{k+1}-x^{\star}\|^{2}
\le &\; \|x_{k}-x^{\star}\|^{2}
-\|x_{k+1}-x_{k}\|^{2} \\
&-2\eta\Bigl[\,f(x_{k+1})-f(x_{k})-\frac{L}{2}\|x_{k+1}-x_{k}\|^{2}
+\langle\nabla f(x_{k}),\,x_{k}-x^{\star}\rangle\Bigr].
\end{aligned}
\tag{14}
\]

\noindent\textbf{[Step 10]} Distribute $-2\eta$ and collect the squared‑norm terms:
\[
\begin{aligned}
\|x_{k+1}-x^{\star}\|^{2}
\le &\; \|x_{k}-x^{\star}\|^{2}
-2\eta\bigl[f(x_{k+1})-f(x_{k})\bigr] \\
&-\Bigl(1-\eta L\Bigr)\|x_{k+1}-x_{k}\|^{2}
-2\eta\langle\nabla f(x_{k}),\,x_{k}-x^{\star}\rangle .
\end{aligned}
\tag{15}
\]

\noindent\textbf{[Step 11]} By convexity of $f$,
\[
f(x_{k})-f(x^{\star})\le\langle\nabla f(x_{k}),\,x_{k}-x^{\star}\rangle .
\tag{16}
\]
Thus $-2\eta\langle\nabla f(x_{k}),\,x_{k}-x^{\star}\rangle\le-2\eta\bigl[f(x_{k})-f(x^{\star})\bigr]$.

\noindent\textbf{[Step 12]} Substitute \eqref{eq:16} into \eqref{eq:15}:
\[
\begin{aligned}
\|x_{k+1}-x^{\star}\|^{2}
\le &\; \|x_{k}-x^{\star}\|^{2}
-2\eta\bigl[f(x_{k+1})-f(x^{\star})\bigr] \\
&-\bigl(1-\eta L\bigr)\|x_{k+1}-x_{k}\|^{2}.
\end{aligned}
\tag{17}
\]

\noindent\textbf{[Step 13]} Add $g(x_{k+1})$ to both sides of the inequality $-2\eta\bigl[f(x_{k+1})-f(x^{\star})\bigr]$ and use the definition of the proximal step:
the optimality condition of \eqref{eq:1} yields
\[
\frac{1}{\eta}(y_{k}-x_{k+1})\in\partial g(x_{k+1}),
\]
which by convexity of $g$ gives
\[
g(x^{\star})\ge g(x_{k+1})+\langle\tfrac{1}{\eta}(y_{k}-x_{k+1}),\,x^{\star}-x_{k+1}\rangle .
\tag{18}
\]
Since $y_{k}=x_{k}-\eta\nabla f(x_{k})$, the inner product in \eqref{eq:18} expands to
\[
\langle\tfrac{1}{\eta}(x_{k}-\eta\nabla f(x_{k})-x_{k+1}),\,x^{\star}-x_{k+1}\rangle
= \frac{1}{\eta}\langle x_{k}-x_{k+1},\,x^{\star}-x_{k+1}\rangle
-\langle\nabla f(x_{k}),\,x^{\star}-x_{k+1}\rangle .
\]
Rearranging \eqref{eq:18} yields
\[
g(x_{k+1})-g(x^{\star})
\le \frac{1}{\eta}\langle x_{k+1}-x_{k},\,x_{k+1}-x^{\star}\rangle
+\langle\nabla f(x_{k}),\,x_{k+1}-x^{\star}\rangle .
\tag{19}
\]

\noindent\textbf{[Step 14]} Combine \eqref{eq:19} with the term $-2\eta\langle\nabla f(x_{k}),\,x_{k+1}-x^{\star}\rangle$ that already appears in \eqref{eq:15}.  After simple algebra we obtain
\[
\begin{aligned}
&-2\eta\bigl[f(x_{k+1})-f(x^{\star})\bigr]
-2\eta\bigl[g(x_{k+1})-g(x^{\star})\bigr] \\
\le &\; -2\eta\bigl[F(x_{k+1})-F(x^{\star})\bigr]
+ \|x_{k+1}-x_{k}\|^{2}
-\|x_{k+1}-x^{\star}\|^{2}
+\|x_{k}-x^{\star}\|^{2}.
\end{aligned}
\tag{20}
\]

\noindent\textbf{[Step 15]} Plug \eqref{eq:20} back into \eqref{eq:15} (recall that $F=f+g$) and simplify.  After cancelling the identical $\|x_{k+1}-x^{\star}\|^{2}$ terms we obtain the fundamental inequality
\[
2\eta\bigl[F(x_{k+1})-F(x^{\star})\bigr]
\le \|x_{k}-x^{\star}\|^{2}-\|x_{k+1}-x^{\star}\|^{2}
-(1-\eta L)\|x_{k+1}-x_{k}\|^{2}.
\tag{21}
\]

\noindent\textbf{[Step 16]} Under Assumption~\ref{ass:stepsize} we have $1-\eta L\ge0$, therefore the last term on the right‑hand side is non‑positive and can be dropped:
\[
2\eta\bigl[F(x_{k+1})-F(x^{\star})\bigr]
\le \|x_{k}-x^{\star}\|^{2}-\|x_{k+1}-x^{\star}\|^{2}.
\tag{22}
\]

\noindent\textbf{[Step 17]} Summing \eqref{eq:22} over $k=0,\dots,K-1$ yields a telescoping series:
\[
2\eta\sum_{k=0}^{K-1}\bigl[F(x_{k+1})-F(x^{\star})\bigr]
\le \|x_{0}-x^{\star}\|^{2}-\|x_{K}-x^{\star}\|^{2}
\le \|x_{0}-x^{\star}\|^{2}.
\tag{23}
\]

\noindent\textbf{[Step 18]} Since $F(x_{k})\ge F(x^{\star})$ for all $k$, each term in the sum is non‑negative; dividing \eqref{eq:23} by $2\eta K$ gives
\[
F(x_{K})-F(x^{\star})\;\le\;\frac{\|x_{0}-x^{\star}\|^{2}}{2\eta K},
\]
which is exactly the bound \eqref{eq:2} claimed in Theorem~\ref{thm:rate}. ∎

\section*{Rate Derivation (With Constants)}
The constant in \eqref{eq:2} is $\frac{1}{2\eta}$.
If we choose the largest admissible stepsize $\eta=1/L$, the bound becomes
\[
F(x_{K})-F(x^{\star})\le\frac{L\,\|x_{0}-x^{\star}\|^{2}}{2K}.
\]
Thus the convergence rate is $O\!\bigl(\tfrac{L\|x_{0}-x^{\star}\|^{2}}{K}\bigr)$.

\section*{Special Cases}
\begin{enumerate}[label=(\alph*)]
\item \textbf{Pure Gradient Descent.}  If $g\equiv0$, $\prox_{\eta g}$ is the identity and the algorithm reduces to $x_{k+1}=x_{k}-\eta\nabla f(x_{k})$.  The same proof gives the classical $O(1/K)$ rate.
\item \textbf{Proximal Point Algorithm.}  If $f\equiv0$, the update becomes $x_{k+1}= \prox_{\eta g}(x_{k})$.  The inequality \eqref{eq:22} simplifies to
\[
\|x_{k+1}-x^{\star}\|^{2}\le\|x_{k}-x^{\star}\|^{2},
\]
hence $F(x_{k})\downarrow F(x^{\star})$ at least as fast as $O(1/K)$.
\item \textbf{Strongly Convex $F$.}  If $F$ is $\mu$‑strongly convex, a slight modification of \eqref{eq:22} yields a linear rate
\[
F(x_{k})-F(x^{\star})\le\bigl(1-\eta\mu\bigr)^{k}\bigl[F(x_{0})-F(x^{\star})\bigr].
\]
\end{enumerate}

\end{document}